```latex
\documentclass[11pt,a4paper]{article}

% --------------------------------------------------
% Packages
% --------------------------------------------------
\usepackage{amsmath}
\usepackage{amssymb}
\usepackage{bm}
\usepackage{booktabs}
\usepackage{longtable}
\usepackage{array}
\usepackage{geometry}
\usepackage{hyperref}

\geometry{
    margin=1in
}

\hypersetup{
    colorlinks=true,
    linkcolor=blue,
    urlcolor=blue
}

% --------------------------------------------------
% Custom commands
% --------------------------------------------------
\newcommand{\E}{\mathbb{E}}
\newcommand{\R}{\mathbb{R}}
\newcommand{\dd}{\mathrm{d}}

% --------------------------------------------------
% Document
% --------------------------------------------------
\begin{document}

\begin{titlepage}
    \centering

    \vspace*{2cm}

    {\Huge\bfseries Mathematical English}\\[0.5cm]
    {\LARGE Physics}\\[1.5cm]

    {\large
    English Vocabulary, Japanese Meaning,\\
    Mathematical Formulas, and Physical Meaning
    }

    \vfill

    {\large
    Physics Reference\\
    \today
    }

\end{titlepage}

\tableofcontents
\newpage

% ==================================================
\section{Classical Mechanics}
% ==================================================

\subsection{Basic Concepts}

\begin{longtable}{p{0.18\textwidth} p{0.18\textwidth} p{0.25\textwidth} p{0.30\textwidth}}
\toprule
\textbf{English} & \textbf{日本語} & \textbf{Formula} & \textbf{Meaning} \\
\midrule
\endhead

Position & 位置 &
$x(t)$ &
The location of an object as a function of time. \\

Displacement & 変位 &
$\Delta x = x_f-x_i$ &
Change in position. \\

Velocity & 速度 &
$v=\dfrac{\dd x}{\dd t}$ &
Rate of change of position. \\

Acceleration & 加速度 &
$a=\dfrac{\dd v}{\dd t}
=\dfrac{\dd^2x}{\dd t^2}$ &
Rate of change of velocity. \\

Mass & 質量 &
$m$ &
Measure of an object's inertia. \\

Momentum & 運動量 &
$\bm{p}=m\bm{v}$ &
Quantity of motion carried by an object. \\

Force & 力 &
$\bm{F}=m\bm{a}$ &
Interaction that changes an object's motion. \\

Work & 仕事 &
$W=\int \bm{F}\cdot\dd\bm{r}$ &
Energy transferred by a force acting through a displacement. \\

Kinetic Energy & 運動エネルギー &
$K=\dfrac{1}{2}mv^2$ &
Energy associated with motion. \\

Potential Energy & ポテンシャルエネルギー &
$U(\bm{r})$ &
Energy associated with position or configuration. \\

Mechanical Energy & 力学的エネルギー &
$E=K+U$ &
Sum of kinetic and potential energy. \\

Power & 仕事率・動力 &
$P=\dfrac{\dd W}{\dd t}$ &
Rate at which work is performed. \\

\bottomrule
\end{longtable}

% --------------------------------------------------
\subsection{Kinematics}

For constant acceleration,

\[
v=v_0+at
\]

\[
x=x_0+v_0t+\frac{1}{2}at^2
\]

\[
v^2=v_0^2+2a(x-x_0)
\]

These equations describe motion without requiring an explicit analysis of the forces causing the motion.

% --------------------------------------------------
\subsection{Newton's Laws}

\textbf{Newton's First Law}

An object remains at rest or moves with constant velocity unless acted upon by a net external force.

\[
\sum \bm{F}=0
\]

\textbf{Newton's Second Law}

\[
\sum \bm{F}=m\bm{a}
\]

The net force determines the acceleration of an object.

\textbf{Newton's Third Law}

\[
\bm{F}_{AB}=-\bm{F}_{BA}
\]

For every force exerted by one object on another, there is an equal and opposite force.

% --------------------------------------------------
\subsection{Momentum Conservation}

For an isolated system,

\[
\bm{p}_{\mathrm{initial}}
=
\bm{p}_{\mathrm{final}}
\]

or

\[
\sum_i m_i\bm{v}_i
=
\text{constant}.
\]

Momentum conservation is particularly important in collision and interaction problems.

% ==================================================
\section{Rotational Motion}
% ==================================================

\begin{longtable}{p{0.20\textwidth} p{0.18\textwidth} p{0.25\textwidth} p{0.27\textwidth}}
\toprule
\textbf{English} & \textbf{日本語} & \textbf{Formula} & \textbf{Meaning} \\
\midrule
\endhead

Angular Position & 角位置 &
$\theta$ &
Position measured by an angle. \\

Angular Velocity & 角速度 &
$\omega=\dfrac{\dd\theta}{\dd t}$ &
Rate of change of angular position. \\

Angular Acceleration & 角加速度 &
$\alpha=\dfrac{\dd\omega}{\dd t}$ &
Rate of change of angular velocity. \\

Torque & トルク・力のモーメント &
$\bm{\tau}=\bm{r}\times\bm{F}$ &
Rotational effect of a force. \\

Moment of Inertia & 慣性モーメント &
$I=\sum_i m_ir_i^2$ &
Rotational analogue of mass. \\

Angular Momentum & 角運動量 &
$\bm{L}=\bm{r}\times\bm{p}$ &
Rotational analogue of linear momentum. \\

Rotational Energy & 回転エネルギー &
$K_{\mathrm{rot}}=\dfrac12I\omega^2$ &
Kinetic energy associated with rotation. \\

\bottomrule
\end{longtable}

The rotational form of Newton's second law is

\[
\bm{\tau}=I\bm{\alpha}.
\]

For a rigid body rotating about a fixed axis,

\[
L=I\omega.
\]

If the net external torque is zero,

\[
L=\text{constant}.
\]

% ==================================================
\section{Gravitation}
% ==================================================

\subsection{Newton's Law of Universal Gravitation}

\[
F=G\frac{m_1m_2}{r^2}
\]

where

\[
G\approx6.674\times10^{-11}
\ \mathrm{N\,m^2/kg^2}.
\]

\textbf{English:} gravitational force

\textbf{日本語:} 万有引力

\textbf{Meaning:} The gravitational force between two masses decreases with the square of the distance between them.

Near the Earth's surface,

\[
F_g=mg.
\]

Thus,

\[
g\approx9.81\ \mathrm{m/s^2}.
\]

% ==================================================
\section{Oscillations}
% ==================================================

\subsection{Simple Harmonic Motion}

Simple harmonic motion is described by

\[
x(t)=A\cos(\omega t+\phi).
\]

The velocity is

\[
v(t)=-A\omega\sin(\omega t+\phi),
\]

and the acceleration is

\[
a(t)=-A\omega^2\cos(\omega t+\phi).
\]

Therefore,

\[
a=-\omega^2x.
\]

This relationship is the defining characteristic of simple harmonic motion.

% --------------------------------------------------
\subsection{Mass-Spring System}

For a spring,

\[
F=-kx.
\]

The angular frequency is

\[
\omega=\sqrt{\frac{k}{m}}.
\]

The period is

\[
T=2\pi\sqrt{\frac{m}{k}}.
\]

The total mechanical energy is

\[
E=\frac12kx^2+\frac12mv^2
=\frac12kA^2.
\]

% --------------------------------------------------
\subsection{Simple Pendulum}

For small oscillations,

\[
T=2\pi\sqrt{\frac{L}{g}}.
\]

The small-angle approximation is

\[
\sin\theta\approx\theta.
\]

% ==================================================
\section{Waves}
% ==================================================

\begin{longtable}{p{0.20\textwidth} p{0.18\textwidth} p{0.25\textwidth} p{0.27\textwidth}}
\toprule
\textbf{English} & \textbf{日本語} & \textbf{Formula} & \textbf{Meaning} \\
\midrule
\endhead

Wavelength & 波長 &
$\lambda$ &
Distance between corresponding points of consecutive waves. \\

Frequency & 振動数・周波数 &
$f$ &
Number of cycles per unit time. \\

Period & 周期 &
$T=\dfrac1f$ &
Time required for one complete cycle. \\

Wave Speed & 波の速さ &
$v=f\lambda$ &
Speed at which the wave propagates. \\

Amplitude & 振幅 &
$A$ &
Maximum displacement from equilibrium. \\

Wave Number & 波数 &
$k=\dfrac{2\pi}{\lambda}$ &
Spatial frequency of a wave. \\

Angular Frequency & 角周波数 &
$\omega=2\pi f$ &
Frequency expressed in radians per second. \\

\bottomrule
\end{longtable}

A one-dimensional traveling wave can be written as

\[
y(x,t)=A\cos(kx-\omega t+\phi).
\]

% ==================================================
\section{Fluid Mechanics}
% ==================================================

\subsection{Density}

\[
\rho=\frac{m}{V}.
\]

\textbf{Density} = 密度

Density represents mass per unit volume.

% --------------------------------------------------
\subsection{Pressure}

\[
P=\frac{F}{A}.
\]

Pressure is force distributed over an area.

For a fluid at rest,

\[
P=P_0+\rho gh.
\]

% --------------------------------------------------
\subsection{Continuity Equation}

For an incompressible steady flow,

\[
A_1v_1=A_2v_2.
\]

This expresses conservation of mass.

% --------------------------------------------------
\subsection{Bernoulli's Equation}

For an ideal incompressible steady flow,

\[
P+\frac12\rho v^2+\rho gh
=
\text{constant}.
\]

The three terms represent pressure energy, kinetic energy per unit volume, and gravitational potential energy per unit volume.

% ==================================================
\section{Thermodynamics}
% ==================================================

\begin{longtable}{p{0.20\textwidth} p{0.18\textwidth} p{0.25\textwidth} p{0.27\textwidth}}
\toprule
\textbf{English} & \textbf{日本語} & \textbf{Formula} & \textbf{Meaning} \\
\midrule
\endhead

Temperature & 温度 &
$T$ &
Measure related to thermal energy and equilibrium. \\

Heat & 熱 &
$Q$ &
Energy transferred because of a temperature difference. \\

Internal Energy & 内部エネルギー &
$U$ &
Energy contained within a system at the microscopic level. \\

Entropy & エントロピー &
$S$ &
State function associated with energy dispersal and the number of accessible microscopic states. \\

Enthalpy & エンタルピー &
$H=U+PV$ &
Useful thermodynamic potential, especially for constant-pressure processes. \\

\bottomrule
\end{longtable}

% --------------------------------------------------
\subsection{First Law of Thermodynamics}

\[
\Delta U=Q-W.
\]

Here, $W$ is work done by the system.

This equation expresses conservation of energy in thermodynamic systems.

% --------------------------------------------------
\subsection{Ideal Gas Law}

\[
PV=nRT.
\]

where

\[
R=8.314\ \mathrm{J/(mol\,K)}.
\]

% --------------------------------------------------
\subsection{Entropy}

For a reversible process,

\[
\dd S=\frac{\delta Q_{\mathrm{rev}}}{T}.
\]

For an isolated system,

\[
\Delta S\geq0.
\]

This is one mathematical expression of the second law of thermodynamics.

% ==================================================
\section{Electromagnetism}
% ==================================================

\subsection{Electric Charge}

\textbf{Electric charge} = 電荷

The elementary charge is

\[
e\approx1.602\times10^{-19}\ \mathrm{C}.
\]

% --------------------------------------------------
\subsection{Coulomb's Law}

\[
F=k_e\frac{|q_1q_2|}{r^2}.
\]

In vector form,

\[
\bm{F}_{12}
=
\frac{1}{4\pi\varepsilon_0}
\frac{q_1q_2}{r^2}\hat{\bm r}.
\]

% --------------------------------------------------
\subsection{Electric Field}

\[
\bm{E}=\frac{\bm{F}}{q}.
\]

For a point charge,

\[
\bm{E}
=
\frac{1}{4\pi\varepsilon_0}
\frac{q}{r^2}\hat{\bm r}.
\]

% --------------------------------------------------
\subsection{Electric Potential}

\[
V=\frac{U}{q}.
\]

For a point charge,

\[
V=
\frac{1}{4\pi\varepsilon_0}
\frac{q}{r}.
\]

The relationship between electric field and potential is

\[
\bm{E}=-\nabla V.
\]

% --------------------------------------------------
\subsection{Ohm's Law}

\[
V=IR.
\]

where

\begin{itemize}
    \item $V$ = voltage
    \item $I$ = electric current
    \item $R$ = resistance
\end{itemize}

Electrical power is

\[
P=VI=I^2R=\frac{V^2}{R}.
\]

% ==================================================
\section{Maxwell's Equations}
% ==================================================

Maxwell's equations provide a compact mathematical description of classical electromagnetism.

\subsection{Gauss's Law for Electricity}

\[
\nabla\cdot\bm{E}
=
\frac{\rho}{\varepsilon_0}.
\]

\subsection{Gauss's Law for Magnetism}

\[
\nabla\cdot\bm{B}=0.
\]

This expresses the absence of isolated magnetic monopoles in classical electromagnetism.

\subsection{Faraday's Law}

\[
\nabla\times\bm{E}
=
-\frac{\partial\bm{B}}{\partial t}.
\]

A changing magnetic field produces an electric field.

\subsection{Amp\`ere--Maxwell Law}

\[
\nabla\times\bm{B}
=
\mu_0\bm{J}
+
\mu_0\varepsilon_0
\frac{\partial\bm{E}}{\partial t}.
\]

A current and a changing electric field produce a magnetic field.

% ==================================================
\section{Special Relativity}
% ==================================================

\subsection{Lorentz Factor}

\[
\gamma=
\frac{1}{\sqrt{1-\frac{v^2}{c^2}}}.
\]

where

\[
c\approx3.00\times10^8\ \mathrm{m/s}.
\]

% --------------------------------------------------
\subsection{Time Dilation}

\[
\Delta t=\gamma\Delta\tau.
\]

Moving clocks are measured to run slower relative to an observer in another inertial frame.

% --------------------------------------------------
\subsection{Length Contraction}

\[
L=\frac{L_0}{\gamma}.
\]

% --------------------------------------------------
\subsection{Mass-Energy Equivalence}

\[
E=mc^2.
\]

More generally, relativistic energy and momentum satisfy

\[
E^2=p^2c^2+m^2c^4.
\]

% ==================================================
\section{Quantum Mechanics}
% ==================================================

\subsection{Planck Relation}

\[
E=hf=\hbar\omega.
\]

where

\[
\hbar=\frac{h}{2\pi}.
\]

% --------------------------------------------------
\subsection{de Broglie Wavelength}

\[
\lambda=\frac{h}{p}.
\]

Particles can exhibit wave-like behavior.

% --------------------------------------------------
\subsection{Heisenberg Uncertainty Principle}

\[
\Delta x\,\Delta p
\geq
\frac{\hbar}{2}.
\]

Position and momentum cannot simultaneously have arbitrarily precise values.

% --------------------------------------------------
\subsection{Schr\"odinger Equation}

The time-dependent Schr\"odinger equation is

\[
i\hbar
\frac{\partial\psi}{\partial t}
=
\hat{H}\psi.
\]

For a non-relativistic particle in a potential $V$,

\[
i\hbar
\frac{\partial\psi}{\partial t}
=
\left[
-\frac{\hbar^2}{2m}\nabla^2
+
V
\right]\psi.
\]

The time-independent equation is

\[
\hat{H}\psi=E\psi.
\]

% --------------------------------------------------
\subsection{Probability Density}

The probability density is

\[
P(x)=|\psi(x)|^2.
\]

The wave function must satisfy the normalization condition

\[
\int |\psi(x)|^2\,\dd x=1.
\]

% ==================================================
\section{Statistical Mechanics}
% ==================================================

\subsection{Boltzmann Distribution}

\[
P(E)\propto e^{-E/(k_BT)}.
\]

where $k_B$ is the Boltzmann constant.

% --------------------------------------------------
\subsection{Boltzmann Entropy}

\[
S=k_B\ln\Omega.
\]

Here, $\Omega$ represents the number of accessible microscopic states.

% --------------------------------------------------
\subsection{Partition Function}

\[
Z=\sum_i e^{-\beta E_i},
\]

where

\[
\beta=\frac{1}{k_BT}.
\]

The partition function is a central quantity in statistical mechanics.

% ==================================================
\section{Lagrangian and Hamiltonian Mechanics}
% ==================================================

\subsection{Lagrangian}

The Lagrangian is defined as

\[
L=T-U.
\]

The equations of motion follow from the Euler--Lagrange equation:

\[
\frac{\dd}{\dd t}
\left(
\frac{\partial L}{\partial\dot q_i}
\right)
-
\frac{\partial L}{\partial q_i}
=0.
\]

This formulation can describe classical mechanics without directly resolving every force component.

% --------------------------------------------------
\subsection{Hamiltonian}

The Hamiltonian is

\[
H=\sum_i p_i\dot q_i-L.
\]

Hamilton's equations are

\[
\dot q_i=\frac{\partial H}{\partial p_i},
\]

\[
\dot p_i=-\frac{\partial H}{\partial q_i}.
\]

In many physical systems, the Hamiltonian corresponds to the total energy.

% ==================================================
\section{Core Physical Constants}
% ==================================================

\begin{longtable}{p{0.25\textwidth} p{0.25\textwidth} p{0.35\textwidth}}
\toprule
\textbf{Constant} & \textbf{Symbol} & \textbf{Approximate Value} \\
\midrule
Speed of light & $c$ & $3.00\times10^8\ \mathrm{m/s}$ \\

Gravitational constant & $G$ & $6.674\times10^{-11}\ \mathrm{N\,m^2/kg^2}$ \\

Planck constant & $h$ & $6.626\times10^{-34}\ \mathrm{J\,s}$ \\

Reduced Planck constant & $\hbar$ & $1.055\times10^{-34}\ \mathrm{J\,s}$ \\

Boltzmann constant & $k_B$ & $1.381\times10^{-23}\ \mathrm{J/K}$ \\

Elementary charge & $e$ & $1.602\times10^{-19}\ \mathrm{C}$ \\

Vacuum permittivity & $\varepsilon_0$ & $8.854\times10^{-12}\ \mathrm{F/m}$ \\

Vacuum permeability & $\mu_0$ & approximately $1.257\times10^{-6}\ \mathrm{H/m}$ \\

\bottomrule
\end{longtable}

% ==================================================
\section{Common Physics Symbols}
% ==================================================

\begin{longtable}{p{0.20\textwidth} p{0.30\textwidth} p{0.35\textwidth}}
\toprule
\textbf{Symbol} & \textbf{English} & \textbf{日本語} \\
\midrule
$x$ & position & 位置 \\
$v$ & velocity & 速度 \\
$a$ & acceleration & 加速度 \\
$m$ & mass & 質量 \\
$\bm{p}$ & momentum & 運動量 \\
$\bm{F}$ & force & 力 \\
$E$ & energy & エネルギー \\
$K$ & kinetic energy & 運動エネルギー \\
$U$ & potential energy & ポテンシャルエネルギー \\
$P$ & power / pressure & 仕事率 / 圧力 \\
$\omega$ & angular frequency & 角周波数 \\
$\lambda$ & wavelength & 波長 \\
$f$ & frequency & 周波数 \\
$\theta$ & angle & 角度 \\
$\rho$ & density / charge density & 密度 / 電荷密度 \\
$\psi$ & wave function & 波動関数 \\
$H$ & Hamiltonian & ハミルトニアン \\
$L$ & Lagrangian / angular momentum & ラグランジアン / 角運動量 \\
\bottomrule
\end{longtable}

% ==================================================
\section{Key Distinctions}
% ==================================================

\subsection{Mass vs. Weight}

\textbf{Mass} is a measure of inertia:

\[
m.
\]

\textbf{Weight} is a gravitational force:

\[
W=mg.
\]

Mass is measured in kilograms, while weight is measured in newtons.

% --------------------------------------------------
\subsection{Energy vs. Power}

Energy measures the amount of energy transferred or stored:

\[
E.
\]

Power measures the rate of energy transfer:

\[
P=\frac{\dd E}{\dd t}.
\]

% --------------------------------------------------
\subsection{Velocity vs. Acceleration}

Velocity describes the rate of change of position:

\[
v=\frac{\dd x}{\dd t}.
\]

Acceleration describes the rate of change of velocity:

\[
a=\frac{\dd v}{\dd t}.
\]

An object can have nonzero velocity while having zero acceleration.

% ==================================================
\section{Physics Learning Map}
% ==================================================

A useful conceptual progression is:

\[
\boxed{
\text{Mechanics}
\rightarrow
\text{Waves}
\rightarrow
\text{Thermodynamics}
\rightarrow
\text{Electromagnetism}
\rightarrow
\text{Relativity}
\rightarrow
\text{Quantum Mechanics}
}
\]

Mathematics provides the language for describing these physical theories.

A more mathematical progression is:

\[
\text{Algebra}
\rightarrow
\text{Calculus}
\rightarrow
\text{Differential Equations}
\rightarrow
\text{Linear Algebra}
\rightarrow
\text{Vector Calculus}
\rightarrow
\text{Advanced Physics}.
\]

% ==================================================
\section{Core Relationships}
% ==================================================

Some of the most important relationships in introductory physics are:

\[
\boxed{F=ma}
\]

\[
\boxed{p=mv}
\]

\[
\boxed{E=K+U}
\]

\[
\boxed{W=\Delta E}
\]

\[
\boxed{E=mc^2}
\]

\[
\boxed{E=hf}
\]

\[
\boxed{\lambda=\frac{h}{p}}
\]

\[
\boxed{
i\hbar\frac{\partial\psi}{\partial t}
=
\hat H\psi
}
\]

These equations illustrate how physical concepts are expressed mathematically.

% ==================================================
\section{Physics and Engineering}
% ==================================================

Physics provides fundamental models of nature, while engineering applies these models to designed systems.

For example,

\[
F=ma
\]

can be used in mechanical system design,

\[
V=IR
\]

in electrical systems,

\[
PV=nRT
\]

in thermodynamic systems, and

\[
i\hbar\frac{\partial\psi}{\partial t}
=
\hat H\psi
\]

in quantum systems.

For computational and engineering applications, physical models are often combined with numerical methods, optimization, simulation, and data analysis.

% ==================================================
\end{document}
```
